\documentclass[11pt,a4paper]{article}
\usepackage[T1]{fontenc}
\usepackage{lmodern,amsmath,amssymb,amsthm,microtype,hyperref}
\usepackage[margin=27mm]{geometry}
\setlength{\emergencystretch}{2em}
\newcommand{\Q}{\mathbf Q}
\newcommand{\Pp}{\mathbf P}
\newtheorem{theorem}{Theorem}
\newtheorem{lemma}[theorem]{Lemma}
\newtheorem{proposition}[theorem]{Proposition}
\newtheorem{corollary}[theorem]{Corollary}
\title{Coupled physical weights and exact rational lifting}
\author{}\date{17 November 2026}
\begin{document}\maketitle
\begin{abstract}
The two physical weights of the general actual-trace
source determine its common slope birationally, before and after
their rational normalization, for every fixed
positive rational comparison scale. The proof uses the two unique
poles and the positive rational contact set. A coprimality certificate
handles the scale at which the poles coincide. This gives an exact
coupled lifting algorithm: a generic compatible rational pair has a
rational slope obtained from a linear polynomial greatest common
divisor. Exceptional image points require the full rational-root
test. Normalized weights can be lifted directly, or through up to
four reciprocal branch pairs.
The two actual-trace square equations and cofinite integer-input
realization remain unproved. No smaller complete defining count is
claimed.
\end{abstract}

\section{The physical pair}

Use the reweighted general source of \cite{Reweighted}, its trace
criterion \cite{Trace}, and the full auxiliary charts of \cite{Auxiliary}.
Fix its rational input and root point tuple, hence its nonzero rational
scalar $\rho$ and positive rational comparison scale
\[
                    a=1+\rho T^2.
\]
The result below applies to every positive rational $a$ and every
$\rho\ne0$, independently of these particular formulas.
Set
\[
 \Theta=\{0,-1/2,-1,1/2399,-32/407\},\qquad
 \Theta_0=(2+\Theta)\cup(4+\Theta),\qquad
 P(t)=\prod_{\theta\in\Theta_0}(t-\theta).
\]
The ten elements of $\Theta_0$ are distinct positive rational numbers.
Since the second contact set is $\Theta_0+4$, the physical pair is
\begin{equation}\label{eq:pair}
 u_0(\lambda)=\rho\frac{P(\lambda)}{(\lambda+1)^{10}},\qquad
 u_1(\lambda)=\rho\frac{P(a\lambda-4)}{(a\lambda+2)^{10}}.
\end{equation}
Both functions are of degree ten. Their denominators are the original
contact denominators, and their unique poles are respectively $-1$
and $-2/a$, each of order ten. There is no cancellation at these
poles because every root of $P$ is positive.

\section{A birationality theorem}

\begin{lemma}\label{lem:two-poles}
Let $f_0,f_1\in\Q(t)$ be nonconstant. Suppose that each has a
unique pole, that these poles $p_0,p_1$ are distinct rational points
of $\Pp^1$, and that all zeros of $f_0$ are distinct rational points.
Let $z$ be a rational coordinate with its zero at $p_0$ and pole at
$p_1$. If all zeros of $f_0$ have positive $z$-coordinate, then
\[
                         \Q(f_0,f_1)=\Q(t).
\]
\end{lemma}
\begin{proof}
By L\"uroth's theorem, the intermediate field is $\Q(h)$ for
some nonconstant $h\in\Q(t)$; a formal reference for precisely
this field statement is \cite{Luroth}. Write $f_i=g_i\circ h$,
and let $d=\deg h$.
Every point above a pole of $g_i$ is a pole of $f_i$. Since $f_i$
has just one pole, $g_i$ has one pole $c_i=h(p_i)$ and its full
fiber is $d p_i$. The two values $c_0,c_1$ are distinct: equality
would make $p_1$ a pole of $f_0$. They are rational because $h$
and $p_i$ are rational.

Apply a rational fractional linear transformation to $h$ sending
$c_0$ to zero and $c_1$ to infinity. The resulting function has
divisor $d p_0-d p_1$, so it equals $c z^d$ with
$c\in\Q^\times$. In particular both $f_i$ are rational
functions of $z^d$.

Choose a zero $t_0$ of $f_0$ and put $z_0=z(t_0)>0$.
Its $z^d$-fiber consists of the $d$ distinct values
$\zeta z_0$, where $\zeta^d=1$, and all give zeros of $f_0$.
If $d\ge3$, some $\zeta$ is nonreal, contradicting the
rationality of every zero. If $d=2$, the second value is $-z_0$,
contradicting positivity of every zero's $z$-coordinate. Thus
$d=1$, proving the claim.
\end{proof}

\begin{theorem}\label{thm:birational}
For every $\rho\in\Q^\times$ and $a\in\Q_{>0}$,
the functions in \eqref{eq:pair} satisfy
\[
                       \Q(u_0,u_1)=\Q(\lambda).
\]
The map from the common slope to its physical weight pair is
therefore birational onto its image curve.
\end{theorem}
\begin{proof}
If $a\ne2$, the poles are distinct. Use
\[
                         z=\frac{\lambda+1}{a\lambda+2}.
\]
At every contact $\theta\in\Theta_0$ both numerator and
denominator are positive. Lemma~\ref{lem:two-poles} applies.
For example, the only possible rational nontrivial deck involution
would be
\[
 \lambda\longmapsto
 -\frac{(a+2)\lambda+4}{2a\lambda+a+2},
\]
which sends every positive contact to a negative number and hence
cannot preserve the contact set.

It remains to handle $a=2$. In the coordinate $s=1/(\lambda+1)$,
division by $\rho$ gives the polynomial pair
\[
 q_0(s)=\prod_{\theta\in\Theta_0}(1-(\theta+1)s),\qquad
 q_1(s)=\prod_{\theta\in\Theta_0}(1-(\theta+6)s/2).
\]
These both have degree ten. If their joint field were $\Q(h)$
with $\deg h>1$, the unique pole of $q_0$ would permit a
fractional linear change of $h$ making $h$ a polynomial, with
both $q_i$ polynomials in $h$. Thus the nonconstant polynomial
$h'$ would divide both $q_0'$ and $q_1'$.

The exact certificate below gives $\gcd(q_0',q_1')=1$ in
$\Q[s]$, a contradiction. This handles the coincident-pole
case as well.
\end{proof}

For the original scalar of the general source, the coincident-pole
case is in fact impossible: $v_{76249}(\rho)$ is odd, whereas
$a=2$ would give $\rho T^2=1$. The proof includes this case so
that the lifting theorem is independent of that particular scalar.

\section{The coincident-pole certificate}

Multiply each $q_i$ by the positive least common multiple of its
coefficient denominators, obtaining integer polynomials $Q_i$.
Modulo $101$, their derivatives, in descending coefficient order,
are
\[
\begin{split}
 f={}&(43,35,20,27,30,29,73,42,14,73),\\
 g={}&(13,58,99,60,62,61,95,88,41,49).
\end{split}
\]
Their degrees are both nine, so reduction preserves both derivative
degrees. The coefficient lists
\[
\begin{split}
 e={}&(3,74,9,88,39,35,15,93,7),\\
 k={}&(60,0,32,85,50,36,38,14,37)
\end{split}
\]
satisfy $ef+kg=1$ in $\mathbf F_{101}[s]$.
Consequently the rational derivatives are coprime. The primary
checker constructs the integer polynomials directly from the contact
set. A separate checker, using only integer lists, differentiates
them modulo $101$ and verifies this Bezout identity by convolution.
The full integer coefficient lists are in the machine-readable
certificate; no floating-point calculation is used.

\section{Exact lifting, including exceptional fibers}

\begin{lemma}[Normalization without lost joint information]\label{lem:normalization}
Let $f_1,\ldots,f_r\in\Q(t)^\times$ be nonconstant and generate
$\Q(t)$ as a field over $\Q$. Suppose that, for every $i$, no
fractional linear transformation of $\Pp^1$ defined over $\Q$
pulls the pole divisor of $f_i$ back to its zero divisor.
Then
\[
 \Q\left(\frac{f_1}{(1+f_1)^2},\ldots,
          \frac{f_r}{(1+f_r)^2}\right)=\Q(t).
\]
\end{lemma}
\begin{proof}
Denote the field on the left by $K$ and put $L=\Q(t)$.
Each $f_i$ satisfies
\[
 w_i Z^2+(2w_i-1)Z+w_i=0,\qquad
 w_i=f_i/(1+f_i)^2\in K.
\]
Its roots are $f_i$ and $f_i^{-1}$, both in $L$.
Thus $L=K(f_1,\ldots,f_r)$ is the splitting field over $K$ of
these separable quadratics. It is a finite Galois extension.
Every $K$-automorphism $\sigma$ of $L$ fixes $\Q$ and sends
each $f_i$ to either $f_i$ or $f_i^{-1}$.

A $\Q$-automorphism of $\Q(t)$ sends $t$ to a fractional linear
function: the degrees of this function and its inverse multiply to
one. The alternative $\sigma(f_i)=f_i^{-1}$ would therefore
identify the indicated pole and zero divisors, contrary to the
hypothesis. Every $\sigma$ fixes every $f_i$, hence fixes $L$.
The finite Galois group is trivial, so $K=L$.
\end{proof}

\begin{corollary}\label{cor:normalized-birational}
For the physical pair \eqref{eq:pair}, put
$w_j=u_j/(1+u_j)^2$. For every positive rational $a$ and nonzero
rational $\rho$,
\[
                         \Q(w_0,w_1)=\Q(\lambda).
\]
\end{corollary}
\begin{proof}
The unnormalized functions jointly generate $\Q(\lambda)$ by
Theorem~\ref{thm:birational}. Each has a single pole of order ten
and ten distinct simple zeros. A fractional linear transformation
preserves the number and multiplicities of points in a divisor,
so it cannot identify these two divisors. Apply
Lemma~\ref{lem:normalization}.
\end{proof}

\begin{proposition}\label{prop:lift}
Fix $\rho\in\Q^\times$, $a\in\Q_{>0}$, and a rational candidate
pair $(v_0,v_1)$. Form
\[
\begin{split}
 F_0(t)&=\rho P(t)-v_0(t+1)^{10},\\
 F_1(t)&=\rho P(at-4)-v_1(at+2)^{10},\\
 G(t)&=\gcd(F_0,F_1)\quad\text{in }\Q[t].
\end{split}
\]
The finite rational slopes producing this pair are precisely the
rational roots $r$ of $G$ with $r\ne-1$ and $ar\ne-2$.
If $G$ is linear, its unique root is rational and gives the unique
finite physical lift, subject to these two pole checks.
If $G$ is constant, there is no finite lift. For $\deg G>1$,
every rational root must be tested; no rationality conclusion follows
solely from the rationality of $(v_0,v_1)$.
\end{proposition}
\begin{proof}
Away from the two contact poles, the equations $F_i(r)=0$ are
exactly $u_i(r)=v_i$. The common finite roots are the roots of
their greatest common divisor. A linear polynomial over $\Q$
has a rational root, and the other assertions are immediate.
\end{proof}

There is also a direct normalized test. Write
\[
 A_0=\rho P(t),\quad B_0=(t+1)^{10},\qquad
 A_1=\rho P(at-4),\quad B_1=(at+2)^{10}.
\]
For a candidate normalized pair $(w_0,w_1)$ form
\begin{equation}\label{eq:normalized-gcd}
 H_j=w_j(A_j+B_j)^2-A_jB_j\quad(j=0,1).
\end{equation}
The finite rational slopes are exactly the rational roots of
$\gcd(H_0,H_1)$ on $B_0B_1(A_0+B_0)(A_1+B_1)\ne0$.
Each $H_j$ has degree at most twenty. The same linear-gcd rule
holds generically by Corollary~\ref{cor:normalized-birational};
higher-degree gcds must again retain every rational-root test.
This method does not assume that either inverse quadratic has
a rational root in advance.

Theorem~\ref{thm:birational} implies that $G$ is linear for a
generic image pair. There is an inverse rational function on a
nonempty open subset of the image, computable by the first nonzero
linear subresultant. For example, if that subresultant is
$A(v_0,v_1)t+B(v_0,v_1)$ with $A\ne0$, its inverse is
$-B/A$. Only finitely many points of the image curve lie outside
such an open subset. Singular image points, branch collisions and
the image of the infinite slope must not be inferred to lift
rationally from this statement. Proposition~\ref{prop:lift}
handles them separately and exactly.

\begin{corollary}\label{cor:normalized}
Given nonzero rational normalized weights $w_0,w_1$, first solve
\[
 w_jv_j^2+(2w_j-1)v_j+w_j=0\qquad(j=0,1).
\]
There are at most four rational branch pairs $(v_0,v_1)$.
For each pair, Proposition~\ref{prop:lift} gives exactly its finite
rational common slopes. The union of these lists is the full
rational lift set for the normalized pair.
\end{corollary}
\begin{proof}
The displayed equation is equivalent to
$w_j=v_j/(1+v_j)^2$. Its discriminant is $1-4w_j$, and its
nonzero roots, when rational, are reciprocal. The value $v_j=-1$
never solves it. At $w_j=1/4$ the two roots coincide at one.
Apply Proposition~\ref{prop:lift} to each rational pair.
\end{proof}

\section{What remains for the twenty-one-variable target}

The theorem concerns the physical scalar pair only. In particular,
the general detector contains the contact denominators
$\delta_0=\lambda+1$ and $\delta_1=a\lambda+2$ in addition
to its weights. They must be recovered from the same lifted slope
before evaluating the actual products and traces. One cannot assign
the traces independently or evaluate them using only two arbitrary
scalar weights.

For each fixed actual elliptic tuple, the exact procedure is now:
recover each candidate common slope by the coupled test, reconstruct
both complete detectors, check
\[
 Y_j^2=V_j^2+w_j(2-V_j),\qquad j=0,1,
\]
and check all original point, difference, scalar, contact, denominator
and whole-character guards. Prepared witnesses also require all
fixed-place, selected-contact and real conditions. No step of the
birationality proof asserts these square values or supplies such
tuples for all sufficiently large positive integers. The complete
general count therefore remains twenty-five under the inherited
arithmetic premises; twenty-one is still a sound source with an
unproved completeness condition.

The checked positive fixtures use $\rho=3/7$, $\lambda=7$ and
$a=3/2,2,3$. They test exact physical lifting, and are not prepared
elliptic witnesses. A pair taking its first weight at slope seven
and its second weight at slope eight, with $a=3$, has constant gcd
and no common finite lift. The same accepted and rejected fixtures
are checked directly after normalization by the degree-twenty test.
This explicitly checks that independent
one-weight solutions cannot be accepted as a coupled solution.

\begin{thebibliography}{9}\raggedright
\bibitem{Reweighted}\emph{Quadratic reweightings of native elliptic detectors},
\nolinkurl{papers/pole_parameters/reweighted_elliptic_detectors_20261116_v53.tex}.
\bibitem{Trace}\emph{A rational conic criterion using the actual detector trace},
\nolinkurl{papers/pole_parameters/trace_conic_criterion_20261117_v57.tex}.
\bibitem{Auxiliary}\emph{Full auxiliary point charts for the general twenty-five-variable source},
\nolinkurl{papers/few_variables25/full_auxiliary25_20261117_v72.tex}.
\bibitem{Luroth}The mathlib community, \emph{L\"uroth's theorem},
\texttt{Mathlib.FieldTheory.RatFunc.Luroth}, theorem
\texttt{RatFunc.Luroth.eq\_adjoin\_generator}; checked 17 November 2026.
\url{https://leanprover-community.github.io/mathlib4_docs/Mathlib/FieldTheory/RatFunc/Luroth.html}.
\end{thebibliography}
\end{document}
